\documentclass[11pt,letterpaper]{article}
\usepackage[T1]{fontenc}
\usepackage[utf8]{inputenc}
\usepackage{amsmath,amsthm,mathtools}
\usepackage{newpxtext,newpxmath}
\usepackage[scaled=0.91]{inconsolata}
\usepackage[margin=1in,headheight=15pt]{geometry}
\usepackage{microtype}
\usepackage{booktabs,longtable,array}
\usepackage{xcolor}
\usepackage{listings}
\usepackage{etoolbox}
\usepackage{needspace}
\usepackage{enumitem}
\usepackage{fancyhdr}
\usepackage{titlesec}
\usepackage{tcolorbox}
\usepackage{hyperref}
\usepackage{bookmark}
\definecolor{ink}{HTML}{243A40}
\definecolor{soft}{HTML}{F3F5F3}
\definecolor{rulegray}{HTML}{C9D1CD}
\hypersetup{colorlinks=true,linkcolor=ink,citecolor=ink,urlcolor=ink,
  pdftitle={Exact Functional Decomposition of Polynomials with Algebraic Coefficients},
  pdfauthor={Prepared for Vladimir Reshetnikov},
  pdfsubject={Characteristic-zero theory and a Wolfram Language implementation}}
\urlstyle{same}
\pagestyle{fancy}
\fancyhf{}
\fancyhead[L]{\small Algebraic polynomial decomposition}
\fancyhead[R]{\small Theory and implementation}
\fancyfoot[C]{\thepage}
\renewcommand{\headrulewidth}{0.3pt}
\titleformat{\section}{\Large\bfseries\color{ink}}{\thesection}{0.65em}{}
\titleformat{\subsection}{\large\bfseries\color{ink}}{\thesubsection}{0.65em}{}
\pretocmd{\section}{\Needspace{6\baselineskip}}{}{}
\pretocmd{\subsection}{\Needspace{4\baselineskip}}{}{}
\makeatletter
\pretocmd{\l@section}{\Needspace{3\baselineskip}}{}{}
\makeatother
\setcounter{tocdepth}{2}
\setlength{\parskip}{0.35em}
\setlength{\parindent}{1.1em}
\setlength{\emergencystretch}{2.5em}
\setlist{nosep,leftmargin=1.6em}
\newtheorem{theorem}{Theorem}[section]
\newtheorem{proposition}[theorem]{Proposition}
\newtheorem{lemma}[theorem]{Lemma}
\newtheorem{corollary}[theorem]{Corollary}
\theoremstyle{definition}
\newtheorem{definition}[theorem]{Definition}
\newtheorem{example}[theorem]{Example}
\theoremstyle{remark}
\newtheorem{remark}[theorem]{Remark}
\newcommand{\Q}{\mathbb Q}
\newcommand{\C}{\mathbb C}
\newcommand{\A}{\overline{\mathbb Q}}
\newcommand{\lc}{\operatorname{lc}}
\newcommand{\ord}{\operatorname{ord}}
\newcommand{\poly}{\operatorname{Pol}}
\newcommand{\wt}[1]{\texttt{#1}}
\DeclareMathOperator{\coeff}{coeff}
\lstdefinelanguage{WL}{
  morekeywords={Module,If,While,Do,Table,Return,Join,Prepend,AppendTo,
  Catch,Throw,True,False,Missing,Failure,Root,RootReduce,AlgebraicNumber,
  Expand,CoefficientList,Get,FileNameJoin,Clear,Total,Reverse,First,Last,
  Length,Quotient,ConstantArray,Select,Divisors,Map,MapIndexed,Fold,
  TestReport,VerificationTest,Options,OptionsPattern,OptionValue},
  sensitive=true,morecomment=[s]{(*}{*)},morestring=[b]"
}
\lstset{language=WL,basicstyle=\ttfamily\footnotesize,
  keywordstyle=\bfseries\color{ink},commentstyle=\color{gray},
  stringstyle=\color{ink},showstringspaces=false,columns=fullflexible,
  keepspaces=true,breaklines=true,breakatwhitespace=false,
  frame=single,rulecolor=\color{rulegray},backgroundcolor=\color{soft},
  framesep=6pt,xleftmargin=6pt,xrightmargin=6pt,
  aboveskip=0.8em,belowskip=0.8em,tabsize=2,
  upquote=true}
\BeforeBeginEnvironment{lstlisting}{\par\noindent\begin{minipage}{\linewidth}}
\AfterEndEnvironment{lstlisting}{\end{minipage}\par}
\tcbset{colback=soft,colframe=rulegray,boxrule=0.5pt,arc=0pt,
  left=9pt,right=9pt,top=7pt,bottom=7pt}

\begin{document}
\begin{titlepage}
\thispagestyle{empty}
\vspace*{0.6in}
{\small\sffamily\bfseries\color{ink} EXACT SYMBOLIC COMPUTATION\par}
\vspace{0.3in}
{\Huge\bfseries\color{ink} Functional Decomposition\par}
\vspace{0.08in}
{\huge\bfseries\color{ink} of Polynomials with\par Algebraic Coefficients\par}
\vspace{0.3in}
{\Large Characteristic-zero theory, exact algorithms,\par and a Wolfram Language package\par}
\vspace{0.45in}
{\large Prepared for Vladimir Reshetnikov\par}
\vspace{0.08in}
{7 September 2026 \quad\textbar\quad Version 1.0.0\par}
\vspace{0.5in}
\begin{tcolorbox}
\textbf{Central result.} Once the degree of an inner component is specified,
its monic, zero-constant representative is uniquely determined by the leading
coefficients of the input. It can be constructed without factoring a
polynomial, solving a nonlinear system, choosing an analytic root branch,
or extending the coefficient field.
\end{tcolorbox}
\vspace{0.25in}
\noindent This article supplies proofs of the reconstruction and descent
theorems, a complete decision procedure, duplicate-free enumeration of complete
decompositions, an audit of the derivative-factor method, exact algebraic
arithmetic contracts, and a documented implementation.
\vfill
\noindent\small\textbf{Validation record.} The independent exact reference
implementation passed 574 checks. The package includes 107 native MUnit tests;
they were not executed because the available Wolfram connector returned
HTTP 404 and no local kernel was available. Native execution is distinguished
from mathematical and reference-code validation throughout the article.
\end{titlepage}

\begin{abstract}
Given $p\in\A[x]$, we seek a composition
$p=f_1\circ\cdots\circ f_r$ in which each nonlinear component is
indecomposable. The coefficients may be represented by radicals, polynomial
\wt{Root} objects, or \wt{AlgebraicNumber} values. We develop an exact,
factorization-free procedure valid mathematically over every field of
characteristic zero. For each proper divisor $d$ of $n=\deg p$, put $m=n/d$.
The unique possible normalized right component is the positive-degree part of
the formal $m$th root of $p/\lc(p)$ at infinity. A triangular recurrence
constructs it, and an $h$-adic expansion decides whether it is actually a
right component. This proves both completeness and descent to the original
coefficient field. Recursive splitting yields a complete chain; enumeration
by an indecomposable rightmost component gives all normalized complete
chains without parenthesization duplicates. The accompanying Wolfram
Language package uses exact coefficient vectors and coefficientwise
\wt{RootReduce}. Its elementary dense implementation takes
$O(n^2\tau(n))$ coefficient-field operations for one complete decomposition,
separately from the cost of algebraic-number arithmetic. Classical
factorization-free decomposition and Ritt theory provide the historical and
structural setting; the article derives the working algorithm directly.
\end{abstract}

\tableofcontents
\newpage

\section{The problem, and what the package promises}
\label{sec:scope}

\subsection{Composition, not multiplication}
The problem is to recognize and recover an identity
\begin{equation}\label{eq:basic}
 p(x)=f(h(x)),\qquad \deg f>1,\quad \deg h>1.
\end{equation}
It is unrelated to writing $p$ as a product of irreducible polynomials.
For example, $x^6+x$ has a visible multiplicative factor $x$, but, as we
prove below, it has no nontrivial functional decomposition. Conversely,
compositional behavior cannot be read off just from the irreducible factors
of $p$.

The motivating Mathematica Stack Exchange question asks whether the
integer-coefficient behavior of \wt{Decompose} can be extended to radicals.
Its answer searches through divisors of $p'$ and uses
$(f\circ h)'=(f'\circ h)h'$. The supplied 2019 comment describes a historical
implementation based on this idea \cite{mse}. We do not infer the behavior
of a current Wolfram kernel from that historical statement. The package
here has its own exact arithmetic and decomposition algorithm, rather than
a wrapper that relies on a particular version of the built-in function.
The official \wt{Decompose} convention also puts the outermost polynomial
first \cite{wl-decompose}; we use that order.

Factorization-free decomposition is classical, notably in the work of
Kozen and Landau \cite{kozen-landau}. The aim here is not to claim a new
decomposition paradigm, but to give a self-contained derivation, an exact
algebraic-coefficient implementation, transparent failure behavior, and
sufficient detail to audit or reimplement it.

\subsection{The mathematical domain and the implemented domain}
In the proofs, $K$ is an arbitrary characteristic-zero field. In the
package, $K$ is the number field generated by the evaluated coefficients:
\[
 K=\Q(a_0,\ldots,a_n)\subseteq\A\subset\C.
\]
The implementation accepts exact algebraic numbers, real or complex.
It does not accept unresolved coefficient parameters such as $a$ in
$x^4+a x^2$, even though the mathematics works over $\Q(a)$.
It does not attempt approximate decomposition of floating-point data.

The result of \wt{AlgebraicDecompose[p,x]} is an outermost-first list
$\{f_1,\ldots,f_r\}$ whose composition equals $p$. When $\deg p\geq2$,
every $f_i$ has degree at least two and is indecomposable. Every component
except the first is monic and has constant coefficient zero.
Constants and linear polynomials use the terminal convention $\{p\}$;
they are not described as nonlinear indecomposables.

The other main operations are fixed-degree decomposition, all normalized
right components, all complete normalized chains, an exact decision report,
exact recomposition, and optional verification of completeness and
normalization. Section~\ref{sec:api} specifies these contracts precisely.

\section{Normalization and coefficient rigidity}
\label{sec:rigidity}

\subsection{Removing the affine ambiguity}
For any affine polynomial $\ell(y)=by+c$, $b\ne0$,
\[
 f\circ h=(f\circ\ell)\circ(\ell^{-1}\circ h).
\]
Consequently a request for literally every decomposition has infinitely
many answers, even over $\Q$. A normalization is essential.

\begin{definition}
A right component $h$ is \emph{normalized} if $\lc(h)=1$ and $h(0)=0$.
A chain is normalized if all its components except possibly the first
are normalized.
\end{definition}

Suppose $p=f\circ g$, with $b=\lc(g)$ and $c=g(0)$. Set
\begin{equation}\label{eq:normalization}
 h(x)=\frac{g(x)-c}{b},\qquad F(y)=f(by+c).
\end{equation}
Then $p=F\circ h$, and $h$ is normalized. The outer polynomial is not in
general monic: its leading coefficient must equal $\lc(p)$.

\subsection{The uniquely determined candidate}
Write
\[
 p(x)=a_nx^n+a_{n-1}x^{n-1}+\cdots+a_0,\qquad a_n\ne0.
\]
If $p=f\circ h$ with $\deg h=d$, then $n=md$, where $m=\deg f$.
The possible nontrivial $d$ are therefore precisely the proper divisors of
$n$ greater than one.

\begin{theorem}[Fixed-degree rigidity]\label{thm:rigidity}
Let $d\mid n$, $1<d<n$, and $m=n/d$. There is at most one pair
$(f,h)$ satisfying $p=f\circ h$, $\deg h=d$, and $h$ normalized.
Moreover, the coefficients of its right component are determined by
$a_n,a_{n-1},\ldots,a_{n-d+1}$ using only field arithmetic and division
by nonzero integers.
\end{theorem}

\begin{proof}
Write
\[
 h=x^d+b_1x^{d-1}+\cdots+b_{d-1}x,
 \qquad f=a_n y^m+u_{m-1}y^{m-1}+\cdots+u_0.
\]
All terms of $f(h)-a_nh^m$ have degree at most $d(m-1)=n-d$.
Hence, for $1\leq k<d$, the coefficient of $x^{n-k}$ in $p$ must be
that in $a_nh^m$. In expanding $h^m$, the term involving the new unknown
$b_k$ comes from selecting $b_kx^{d-k}$ in exactly one factor and the
leading term in all the others. Any other contribution involves only
$b_1,\ldots,b_{k-1}$. Thus
\begin{equation}\label{eq:triangular}
 \frac{a_{n-k}}{a_n}=m b_k+
 \Phi_k(b_1,\ldots,b_{k-1})
\end{equation}
for a polynomial $\Phi_k$ with integer coefficients. Since $m\ne0$ in
$K$, these equations determine the $b_k$ successively.

Once $h$ is known, $f$ is unique because substitution by a nonconstant
polynomial is injective: for nonzero $v$, one has
$\deg(v\circ h)=d\deg v$, so $v\circ h\ne0$.
\end{proof}

The theorem concerns \emph{candidates}: solving the top-coefficient
constraints does not prove that a decomposition exists. The lower
coefficients provide an exact consistency test. This distinction is
central to the implementation.

\subsection{Descent: algebraic extensions do not help}

\begin{theorem}[Descent of normalized components]\label{thm:descent}
Let $K\subseteq L$ be fields of characteristic zero and $p\in K[x]$.
If $p$ has a nontrivial decomposition in $L[x]$, its normalized right
component and corresponding outer component belong to $K[x]$.
\end{theorem}

\begin{proof}
Normalize the given decomposition in $L[x]$ by
\eqref{eq:normalization}. Equations~\eqref{eq:triangular} then construct
all coefficients of $h$ from coefficients of $p$ by arithmetic inside
$K$. Thus $h\in K[x]$.

Recover $f$ by repeatedly subtracting a scalar multiple of the appropriate
power of $h$ from $p$, starting at the leading term. Since $h$ is monic,
these scalars and every intermediate residual lie in $K[x]$. Equivalently,
use the monic division algorithm proved in Section~\ref{sec:hadic}.
It follows that all coefficients of $f$ lie in $K$.
\end{proof}

\begin{corollary}\label{cor:absolute}
A polynomial in $K[x]$ is indecomposable over $K$ if and only if it is
indecomposable over any extension field of characteristic zero, including
an algebraic closure of $K$.
\end{corollary}

This is markedly different from multiplicative irreducibility. For
functional decomposition, there is no need to guess which additional
radicals to put into a factorization extension. The field generated by
the input coefficients already contains the normalized answer. A result
may display a different algebraic representation, but it represents an
element of that same field.

\begin{corollary}[Equivariance under embeddings]
Let $\sigma:K\hookrightarrow L$ be a field embedding. Whenever
$(f,h)$ is the normalized decomposition of $p$ at right degree $d$,
$(\sigma(f),\sigma(h))$ is the normalized decomposition of $\sigma(p)$
at degree $d$. A fixed degree has a decomposition before applying
$\sigma$ if and only if it has one afterwards.
\end{corollary}

\begin{proof}
The candidate equations and all consistency equations use rational field
operations. They commute with $\sigma$, which is injective and therefore
preserves nonzero obstructions.
\end{proof}

This explains why selected conjugate \wt{Root} values can be processed
uniformly. It does not license ignoring the selected conjugate: different
embeddings generally give different numerical coefficients.

\section{Formal roots at infinity and an explicit recurrence}
\label{sec:recurrence}

\subsection{A root whose constant term is prescribed}
Put $t=1/x$ and define
\begin{equation}\label{eq:C}
 C(t)=\frac{t^n p(1/t)}{a_n}
     =1+c_1t+c_2t^2+\cdots+c_nt^n,
 \qquad c_j=\frac{a_{n-j}}{a_n}.
\end{equation}
There is a unique formal power series
\[
 B(t)=1+b_1t+b_2t^2+\cdots\in K[[t]]
 \quad\text{such that}\quad B(t)^m=C(t).
\]
For example, its coefficients can be obtained recursively by comparing
powers of $t$: the coefficient of the new $b_k$ in $B^m$ is $m$.
The notation $C^{1/m}$ refers to this \emph{formal} root with constant
term one.

No analytic branch cut is involved. In particular, we never extract an
$m$th root of $a_n$. Dividing by $a_n$ first is what keeps the
construction inside $K$, even when the leading coefficient is a complex
algebraic number whose $m$th roots lie outside $K$.

\begin{proposition}[Positive-degree part formula]\label{prop:formal}
For a given $d\mid n$, $1<d<n$, the unique possible normalized right
component is
\begin{equation}\label{eq:h}
 h_d(x)=x^d+b_1x^{d-1}+\cdots+b_{d-1}x,
 \qquad B=C^{1/m},\quad m=n/d.
\end{equation}
Only $C\bmod t^d$ is required.
\end{proposition}

\begin{proof}
The relation $B^m=C\pmod{t^d}$ gives exactly the top-coefficient
conditions~\eqref{eq:triangular}, after multiplying by $x^n$ and
putting $t=1/x$. Uniqueness from Theorem~\ref{thm:rigidity} identifies
the coefficients.
\end{proof}

\subsection{Derivation of the recurrence}
Differentiating the formal identity $B^m=C$ gives
\[
 mCB'=C'B.
\]
At coefficient $t^{k-1}$, the left side is
$m\sum_{j=0}^{k-1}(k-j)c_jb_{k-j}$ and the right side is
$\sum_{j=1}^{k}j c_jb_{k-j}$. Taking out the $j=0$ term on the
left yields
\begin{equation}\boxed{\displaystyle
 b_0=1,\qquad
 b_k=\frac{1}{mk}\sum_{j=1}^{k}
       \bigl((m+1)j-mk\bigr)c_jb_{k-j}
 \quad(1\leq k<d).}\label{eq:recurrence}
\end{equation}
Every term on the right involves an earlier $b$ coefficient. The cost
of computing the first $d-1$ coefficients is $O(d^2)$ field operations.
The first few are
\begin{align}
 b_1&=\frac{c_1}{m},\label{eq:b1}\\
 b_2&=\frac{c_2}{m}-\frac{m-1}{2m^2}c_1^2,\label{eq:b2}\\
 b_3&=\frac{c_3}{m}-\frac{m-1}{m^2}c_1c_2
       +\frac{(m-1)(2m-1)}{6m^3}c_1^3.\label{eq:b3}
\end{align}
These expressions are also obtained by the formal binomial series
$(1+U)^{1/m}$. The independent reference tests compare the recurrence
with that separately implemented series operation.

\subsection{Relation to the usual approximate root}
In the Laurent series field $K((x^{-1}))$, set
\[
 S(x)=x^d B(x^{-1})=(p(x)/a_n)^{1/m},
\]
where the displayed equality has the formal meaning just specified.
Let $q$ be its full polynomial part, including the constant term. Then
$S-q=O(x^{-1})$, and therefore
\begin{equation}\label{eq:approxroot}
 \deg\bigl(p/a_n-q^m\bigr)<n-d.
\end{equation}
Indeed, $S^m-q^m=(S-q)(S^{m-1}+S^{m-2}q+\cdots+q^{m-1})$
has degree at most $-1+d(m-1)$ at infinity.
This $q$ is the monic approximate $m$th root in the usual polynomial
sense. Our candidate is $h_d=q-q(0)$, so computing $q(0)$ is unnecessary.

If $p=f(h)$ actually holds, expansion in the parameter $1/h$ gives
\begin{equation}\label{eq:formalactual}
 S=h+\frac{u_{m-1}}{m a_n}+O(h^{-1})
   =h+\frac{u_{m-1}}{m a_n}+O(x^{-d}).
\end{equation}
Thus subtracting the constant polynomial part is precisely the
normalization $h(0)=0$, rather than an arbitrary alteration of the
candidate. The upper coefficients determine $h$, and the lower coefficients
must make the formal root have the special shape in
\eqref{eq:formalactual}.

\subsection{Why the recurrence is preferable to symbolic root expansion}
One could try to ask a CAS to expand a radical expression near infinity.
That brings together analytic branch conventions, possible sign cases,
and simplification of symbolic limits, none of which is needed here.
Equation~\eqref{eq:recurrence} only performs exact operations on constants.
It uses no \wt{Series} call and no expression such as
\wt{PowerExpand[(p/a)\string^(1/m)]}. A formal construction with leading
coefficient one is both more predictable and easier to verify.

\section{The exact composition test: expansion in powers of the candidate}
\label{sec:hadic}

\subsection{The \texorpdfstring{$h$}{h}-adic expansion}
Fix any monic polynomial $h$ of degree $d\geq1$. Repeated polynomial
division produces
\begin{align*}
 p&=q_1h+r_0,\\
 q_1&=q_2h+r_1,\\
 &\hspace{0.3em}\vdots\\
 q_s&=r_s,
\end{align*}
where $\deg r_j<d$ and the final quotient is zero after one more division.
Consequently
\begin{equation}\label{eq:hadic}
 p=r_0+r_1h+\cdots+r_sh^s,\qquad \deg r_j<d.
\end{equation}
These $r_j$ are the polynomial analogues of digits in a positional
numeral system.

\begin{lemma}[Uniqueness of $h$-adic digits]\label{lem:digits}
The representation~\eqref{eq:hadic} with $\deg r_j<d$ is unique.
\end{lemma}

\begin{proof}
In a nonzero sum $\sum_jv_jh^j$ with $\deg v_j<d$, choose the largest
$j$ for which $v_j\ne0$. Its degree lies in $[jd,(j+1)d-1]$; every
smaller-index term has degree below $jd$. Therefore its leading term
cannot cancel. Applying this to the difference of two representations
proves uniqueness.
\end{proof}

\begin{theorem}[Composition membership criterion]\label{thm:membership}
For monic $h$, a polynomial $p$ belongs to $K[h]$ if and only if every
$h$-adic digit $r_j$ is constant. In that case the outer component is
\[
 f(y)=\sum_j r_jy^j.
\]
\end{theorem}

\begin{proof}
Constant digits give the displayed composition immediately. Conversely,
if $p=f(h)=\sum_j u_jh^j$ with $u_j\in K$, these constants already form
a valid $h$-adic expansion. Lemma~\ref{lem:digits} forces $r_j=u_j$.
\end{proof}

The first nonconstant remainder is therefore a conclusive obstruction,
not merely a failed attempt to solve for an outer polynomial. Once the
candidate degree is fixed, Theorem~\ref{thm:rigidity} excludes every
alternative normalized right component of that degree.

\subsection{A directly checkable decision record}
For an unsuccessful degree $d$, the report stores $h_d$, the index $j$
of the first nonconstant digit, that remainder, and the previous constant
digits $u_0,\ldots,u_{j-1}$. A checker can recover
\[
 q_j=\frac{p-\sum_{i<j}u_i h_d^i}{h_d^j}
\]
by exact polynomial division and verify that the remainder of $q_j$
modulo $h_d$ is the reported nonconstant polynomial.
It can also reconstruct $h_d$ from the recurrence to check that the
obstruction applies to the only normalized candidate.

For a successful degree, exact recomposition $f(h_d)-p=0$ verifies the
pair. To verify indecomposability of $p$, check a negative record for
every proper degree divisor. The package's report is an exact decision
record suitable for independent checking, but not a proof object already
accepted by a proof-assistant kernel.

\subsection{A leading-term alternative}
An equivalent recovery algorithm subtracts multiples of $h^j$ from the
leading term downward. If a nonzero residual has degree not divisible
by $d$, it cannot be a polynomial in $h$ and the trial fails. Otherwise,
subtract its leading coefficient times the corresponding power of $h$
and continue. This also proves uniqueness of the outer polynomial.

Repeated monic division was chosen for the implementation because it
never constructs all powers of $h$ at once, has a small and explicit
coefficient-domain interface, and directly returns the obstruction digits.
It does not invoke \wt{Solve}: the apparently nonlinear decomposition
problem has become one deterministic candidate computation followed by
polynomial division.

\section{Decision, complete decomposition, and exhaustive enumeration}
\label{sec:algorithms}

\subsection{A complete fixed-degree decision procedure}
For a specified degree $d$, the algorithm is:
\begin{enumerate}
\item Validate exact coefficients, compute the actual degree $n$, and require
$1<d<n$ with $d\mid n$.
\item Put $m=n/d$, form $c_1,\ldots,c_{d-1}$, and compute
$b_1,\ldots,b_{d-1}$ using~\eqref{eq:recurrence}.
\item Form $h_d$ as in~\eqref{eq:h}, compute its $h_d$-adic digits,
and reject at the first nonconstant digit.
\item If all digits are constant, form $f$ from them and verify
$f\circ h_d=p$ exactly before returning the pair.
\end{enumerate}

\begin{theorem}[Soundness, completeness, and termination]\label{thm:fixed}
The fixed-degree procedure terminates. It succeeds exactly when $p$ has
a right component of degree $d$, even if decompositions over extension
fields are allowed. On success, it returns the unique normalized pair.
\end{theorem}

\begin{proof}
The recurrence has $d-1$ finite steps; each monic division and the
sequence of divisions terminate by decreasing degree. On success,
Theorem~\ref{thm:membership} proves equality. If a decomposition exists,
Theorems~\ref{thm:rigidity} and~\ref{thm:descent} identify its normalized
right component with the constructed candidate. Its digits are then
constant by Theorem~\ref{thm:membership}, so the procedure cannot reject.
\end{proof}

Testing all proper divisors of $n$ therefore decides whether $p$ has any
nontrivial decomposition. A polynomial of prime degree needs no trial.
A polynomial of composite degree must still be tested: composite degree
is necessary for decomposition, not sufficient.

\subsection{One complete chain}
To construct one chain, test the possible right degrees in a prescribed
order. At the first success $p=f\circ h$, recursively decompose both
$f$ and $h$, then concatenate the two outermost-first lists. If every
degree fails, return $\{p\}$.

\begin{theorem}[Completeness of the recursive chain]\label{thm:chain}
For $\deg p\geq2$, this recursion terminates and returns a normalized
complete chain. The result is valid for either increasing or decreasing
trial-degree order.
\end{theorem}

\begin{proof}
Both child degrees are strictly smaller than their parent's. Hence the
recursion terminates. Inductively, concatenation preserves composition,
and every returned leaf has no nontrivial right degree by
Theorem~\ref{thm:fixed}.

Normalization also propagates. If the parent is monic and zero at zero,
then in a split with normalized right component, the outer component is
monic because leading coefficients multiply in the usual compositional
way, and its value at zero is the parent's value at zero. Thus both
children are normalized. Starting at the original input, only the
outermost branch can retain a nonunit leading coefficient or a nonzero
constant term.
\end{proof}

Trying the smallest successful right degree also implies that the
selected right component is indecomposable: a proper right component
of it would give a smaller successful degree for $p$. Nevertheless, the
implementation recursively processes both sides. Its correctness does
not depend on that optimization or on a hidden property of trial order.

\subsection{Why arbitrary binary splitting duplicates answers}
Suppose $p=A\circ B\circ C$. A recursion that enumerates all binary
splits can reach the same final chain from $(A\circ B)\circ C$ and from
$A\circ(B\circ C)$. Exact algebraic equality makes post hoc expression
deduplication more cumbersome than necessary. The duplication is caused
by tree parenthesization, not by a genuine second decomposition.

The enumerator instead chooses the \emph{rightmost indecomposable}
component. For each normalized pair $p=f\circ h$, retain it only if
$h$ is indecomposable; enumerate complete chains of $f$ and append $h$
to each. If there are no nontrivial pairs, emit the singleton chain.

\begin{theorem}[Enumeration without duplicates]\label{thm:all}
The preceding procedure enumerates exactly all normalized complete
chains, with each affine-equivalence class represented once. No two
outputs have the same ordered degree sequence.
\end{theorem}

\begin{proof}
A complete chain has a rightmost indecomposable component $h$. At its
degree, the fixed-degree test returns precisely that normalized $h$ and
its unique outer component $f$. The recursion on $f$ supplies the prefix
of the chain, by induction on degree. Thus every complete chain appears.

For uniqueness, the last degree determines $h$ and $f$ uniquely. The
same argument applied to the prefix determines every earlier component.
Induction therefore shows that a degree sequence determines at most one
normalized chain and that no recursive branch duplicates another.
\end{proof}

The code memoizes normalized right-pair lists locally during an
\wt{AlgebraicDecomposeAll} call. It does not install an unbounded global
cache of user polynomials.

\subsection{Finite output caps and honest completeness flags}
A finite \wt{MaxDecompositions} value $M$ limits the returned list, not
the mathematical search space. Emitting exactly $M$ answers is not enough
to know whether enumeration was exhaustive. The implementation searches
until either it is exhausted or it has found $M+1$ distinct chains.
Only the latter case sets \wt{EnumerationComplete} to \wt{False} and
discards the extra chain.

Consequently a polynomial with exactly $M$ chains is correctly reported
as fully enumerated with cap $M$. No count of omitted chains is claimed.
The option is not a time or memory limit: candidate construction and
memoized right-pair tests can require substantial work even before the
first output is emitted.

\subsection{Recovering nonnormalized decompositions}
For a normalized chain $F_1,\ldots,F_r$, choose affine polynomials
$\ell_1,\ldots,\ell_{r-1}$. Every chain in its affine-equivalence class
has the form
\begin{align*}
 &F_1\circ\ell_1,\\
 &\ell_1^{-1}\circ F_2\circ\ell_2,\\
 &\hspace{1.4em}\ldots,\\
 &\ell_{r-2}^{-1}\circ F_{r-1}\circ\ell_{r-1},\\
 &\ell_{r-1}^{-1}\circ F_r.
\end{align*}
The adjacent affine maps cancel in composition. These infinitely many
choices are deliberately not enumerated. The returned normalized chains
are a finite, unambiguous set of representatives.

\section{Structural consequences and the role of Ritt theory}
\label{sec:ritt}

\subsection{What is invariant, and what is not}
Ritt's first theorem connects complete decompositions by elementary
changes involving adjacent factors. In characteristic zero, all complete
chains of the same polynomial have the same length and the same
\emph{multiset} of component degrees, though their degree order can differ.
A modern treatment gives the degree statement as Corollary 2.12
\cite{zieve-mueller}. Descent allows the degree invariant to be applied to
the coefficient field used here, not just to $\C$.

This structural theorem is not used as a substitute for an algorithmic
proof. Theorems~\ref{thm:fixed}--\ref{thm:all} establish the decision and
enumeration procedures directly.

\begin{corollary}[Minimizing the largest component degree]\label{cor:minmax}
Among all decompositions of a fixed nonlinear polynomial into nonlinear
polynomial components, every complete decomposition attains the least
possible value of the largest component degree.
\end{corollary}

\begin{proof}
Refining a component of degree $ab$, with $a,b>1$, replaces it by
components of degrees $a$ and $b$, both smaller than $ab$. Thus refinement
to a complete chain cannot increase the maximum degree. Ritt's degree
invariant makes that final maximum independent of the selected complete
chain. No original chain can have a smaller maximum than its own complete
refinement.
\end{proof}

This does not say that all indecomposable components have prime degree.
For the nested example in Section~\ref{sec:nested}, the invariant
multiset is $\{2,3,4\}$; the indecomposable quartic cannot be replaced by
quadratics in another complete chain.

\subsection{Explicit collisions}
Two useful families of identities can be checked directly:
\begin{align}
 T_m\circ T_n&=T_n\circ T_m,\label{eq:cheb}\\
 x^n\circ\bigl(x^sR(x^n)\bigr)
 &=\bigl(x^sR(x)^n\bigr)\circ x^n.\label{eq:powercollision}
\end{align}
Here $T_k$ is the Chebyshev polynomial. The first identity follows by
putting $x=\cos\theta$, using $T_k(\cos\theta)=\cos(k\theta)$,
and then using equality of polynomials on infinitely many points.
Both sides of the second identity equal $x^{ns}R(x^n)^n$.

Under the appropriate coprime degree conditions, these Chebyshev and
power-type identities, together with affine changes, describe the
nontrivial collisions in Ritt's second theorem
\cite[Theorem 2.17]{zieve-mueller}. The package need not recognize these
special forms: it finds the corresponding chains by the same fixed-degree
tests as for a generic polynomial.

\subsection{Monomials give exact enumeration benchmarks}
For $p=x^N$, every candidate is $h_d=x^d$, and every proper divisor $d$
succeeds with outer polynomial $x^{N/d}$. A complete chain consists
precisely of monomials whose degrees are the prime factors of $N$, in
some order. If
\[
 N=\prod_{i=1}^s p_i^{e_i},\qquad r=\sum_i e_i,
\]
then the number of normalized complete chains is
\begin{equation}\label{eq:monocount}
 \frac{r!}{e_1!\cdots e_s!}.
\end{equation}
This follows by counting distinct permutations of the degree multiset;
Theorem~\ref{thm:all} proves that each sequence determines one chain.
These counts give a test oracle independent of a CAS decomposition
routine. For $N=6,12,16,30,36,72$, the counts are respectively
$2,3,1,6,6,10$.

\subsection{Generic indecomposability and exactness}
For a fixed $n$ and a fixed proper divisor $d$, the coefficients of
$h_d$ are rational functions of $a_0,\ldots,a_n$ with denominators
nonzero on the chart $a_n\ne0$. The $h_d$-adic remainders have the same
property. The condition that their nonconstant coefficients vanish can
therefore be expressed by finitely many polynomial equations after
clearing those denominators.

Thus decomposability at a fixed degree is a Zariski-closed condition
inside the degree-$n$ coefficient chart. Taking the finite union over
$d$ preserves closedness. This union is proper: $x^n+x$ lies outside it
for every $n\geq2$. For a nontrivial $d$, its candidate is $x^d$ and the
first remainder is $x$, which is nonconstant; if there is no such $d$,
indecomposability follows from prime degree.

It follows directly that indecomposability is a nonempty Zariski-open
condition in that chart. Exact compositional identities can disappear
under arbitrarily small generic perturbations. This is why replacing
exact algebraic coefficients by decimal approximations changes the
problem, rather than merely slowing the same decision procedure.

\section{Exact algebraic-number arithmetic}
\label{sec:arithmetic}

\subsection{Values, representations, and branch choices}
A coefficient such as $\sqrt2$, a selected root of $z^5-z-1$, or
$\sqrt2+i$ denotes an exact number. The package must preserve that
number, not just its minimal polynomial. Distinct conjugates satisfy the
same minimal polynomial but can give distinct input polynomials.

Radicals have their usual Wolfram Language value semantics. For
example, a principal complex cube root is not silently replaced by a
real cube root. The implementation never calls \wt{PowerExpand} and
never rewrites radicals by identities that are valid only under unstated
sign assumptions.

The package is concerned with \emph{constant algebraic} \wt{Root}
objects. A \wt{Root} expression depending on the polynomial variable is
not an algebraic coefficient. A numerical or transcendental root is also
outside the domain. Membership in \wt{Algebraics}, together with the
representation checks described below, is required \cite{wl-algebraics}.

\subsection{Canonical reduction and exact zero testing}
The arithmetic layer applies \wt{RootReduce} separately to coefficients,
not to an entire polynomial in the indeterminate $x$. Wolfram documents
\wt{RootReduce} as reducing algebraic combinations of exact algebraic
numbers to a canonical algebraic representation; it can simplify a result
to a rational number or a simpler form where appropriate
\cite{wl-rootreduce}.

An integer or rational coefficient takes a fast path and is left alone.
Other supported coefficients are reduced exactly. Equality testing is
then implemented by reducing a difference and comparing it with the
exact integer zero. No tolerance, \wt{Chop}, \wt{PossibleZeroQ}, numerical
sampling, or floating-point norm is used.

It is not necessary for two equal input coefficients to have literally
identical printed forms. What matters is that their reduced difference
is exactly zero. Leading cancellations must be resolved before computing
the degree. For example, if $\alpha^5-\alpha-1=0$, a displayed term
$(\alpha^5-\alpha-1)x^{20}$ contributes nothing to the actual degree.
The package reduces the coefficient vector and trims exact leading zeros.

\subsection{Input validation is a separate phase}
The input expression is expanded and checked with \wt{PolynomialQ} in
an unassigned, nonnumeric symbol. Its coefficient list is then obtained.
Any surviving machine-precision or arbitrary-precision \wt{Real}
coefficient is rejected, before any attempt to reinterpret it as an
exact number.

Each reduced coefficient must pass a closed algebraic-expression grammar:
integers, rationals, exact complex numbers with algebraic parts, suitable
\wt{Root} or \wt{AlgebraicNumber} values, and algebraic sums, products,
or rational powers. Numeric status and exact algebraic-domain membership
are checked as well. A coefficient that cannot be certified in this
supported representation is rejected with a \wt{Failure} object.

This is intentionally stricter than the abstract characteristic-zero
algorithm. For example, the theory permits a rational function field,
but this implementation does not maintain assumptions about a symbolic
parameter or split into specialization cases. A failure to certify an
input is not reported as mathematical indecomposability.

As usual for an ordinary Wolfram function without holding attributes,
arguments evaluate before the function receives them. These checks apply
to the evaluated coefficient values, not to the historical origin of
an expression before evaluation. The variable should be cleared before
calling the API.

\subsection{Why there is no \texttt{Extension} option}
Factoring a derivative needs a chosen coefficient field in which to
factor. The present procedure only needs field arithmetic, so the
normalized answer is already in the field generated by the input.
Theorem~\ref{thm:descent} explains why an extension parameter would not
add missing decompositions in characteristic zero.

The implementation does not explicitly construct one common primitive
element. \wt{RootReduce} handles exact intermediate coefficient values.
This favors a short, auditable arithmetic boundary, at the cost of
potentially repeated algebraic-number conversion work.

An alternative backend can choose $\theta$ with
$K=\Q(\theta)$ and a minimal polynomial $M$, and represent every
coefficient as a residue class in $\Q[z]/(M)$. Addition is vector
addition; multiplication is polynomial multiplication modulo $M$; a
nonzero inverse is obtained by the extended Euclidean algorithm.
The chosen complex embedding of $\theta$ must be retained when converting
results back to Wolfram expressions. The arithmetic itself then takes
place in a single exact quotient field.

The official \wt{ToNumberField} documentation provides common-field
conversion for lists. It distinguishes \wt{Automatic}, which need not
find the smallest field, from the \wt{All} choice that asks for the
smallest common field \cite{wl-tonumberfield}. Such a backend would be
a performance extension to this release, not a prerequisite for its
mathematical completeness. No explicit primitive-element backend is
claimed to be implemented here.

\clearpage
\section{Implementation architecture and complexity}
\label{sec:implementation}

\subsection{Coefficient vectors and invariants}
Internally a polynomial is an ascending coefficient vector:
\[
 [a_0,a_1,\ldots,a_n]\quad\longleftrightarrow\quad
 \sum_{i=0}^n a_i x^i.
\]
The zero polynomial is exactly \wt{\{0\}}, not an empty list.
Vectors have no trailing zero entries except this zero representation.
A normalized degree-$d$ right component has first entry zero and last
entry one. Every arithmetic result is reduced coefficientwise and trimmed.

The trusted arithmetic boundary consists of exact coefficient reduction
and ordinary exact scalar arithmetic. The package itself implements dense
vector addition, multiplication, composition by Horner evaluation, and
monic polynomial division. It does not delegate the decomposition to
\wt{Decompose}, factorization to \wt{Factor}, or coefficient recovery to
\wt{Solve}.

The source is organized in four layers: input and arithmetic checks;
coefficient-vector operations; candidate construction and $h$-adic
recognition; and public decision, recursion, enumeration, and reporting
functions. Failures use a private tagged \wt{Catch}/\wt{Throw} channel,
so an invalid coefficient or arithmetic problem is not confused with a
negative decomposition result.

\subsection{The candidate routine}
The following is the actual core from the supplied package. In this
listing, \wt{red} is the exact coefficient reducer and vectors are
ascending.
\begin{lstlisting}
rightCandidate[p_List, d_Integer] := Module[
  {n = Length[p] - 1, m, c, b, k, j},
  m = Quotient[n, d];
  c = Table[red[p[[n - j + 1]]/Last[p]],
    {j, 0, d - 1}];
  b = ConstantArray[0, d]; b[[1]] = 1;
  Do[
    b[[k + 1]] = red[
      Total[Table[
        ((m + 1) j - m k) c[[j + 1]] b[[k - j + 1]],
        {j, 1, k}]]/(m k)],
    {k, 1, d - 1}];
  Prepend[Reverse[b], 0]
];
\end{lstlisting}
The reversal converts $[1,b_1,\ldots,b_{d-1}]$ to the ascending
coefficients of the positive-degree terms; prepending zero imposes the
constant normalization. All calls to this internal routine are guarded
by the proper-divisor test.

\subsection{Monic division and the membership routine}
For a monic divisor $h$ of degree $d$, descending long division reads
the current leading coefficient $s$, writes it into the quotient, and
subtracts $s x^{k-d}h$. The leading coefficient of $h$ is one, so no
additional inversion is required. Each changed coefficient is reduced
immediately, preserving exact zero and degree invariants.

The membership routine stops at its first nonconstant digit:
\begin{lstlisting}
outerData[p_List, h_List] := Module[
  {q = p, digits = {}, qr, j = 0},
  While[Length[q] > 1,
    qr = vDivideMonic[q, h];
    If[Length[qr[[2]]] > 1,
      Return[<|"Success" -> False, "DigitIndex" -> j,
        "RemainderVector" -> qr[[2]],
        "PriorDigits" -> digits|>]];
    AppendTo[digits, First[qr[[2]]]];
    q = qr[[1]]; j++];
  AppendTo[digits, First[q]];
  <|"Success" -> True, "OuterVector" -> trim[digits]|>
];
\end{lstlisting}
A vector of length one represents a constant, including zero. Successful
pairs are recomposed exactly by the caller. Complete chains and enumerated
chains are also recomposed before being returned. These redundant identity
checks can expose arithmetic or indexing defects, though recomposition
alone does not prove completeness of a decomposition.

\subsection{Cost of one trial}
Count operations in $K$, treating each exact coefficient operation as
one unit. Candidate construction takes $O(d^2)$ such operations.
A monic division of a degree-$N$ polynomial by a degree-$d$ polynomial
costs $O((N-d+1)d)$. In the $h$-adic expansion, the quotient degrees
decrease by $d$, so the total is bounded by
\[
 O\left(d\sum_{j=1}^{m}(n-jd+1)\right)=O(n^2).
\]
Horner recomposition of an outer polynomial of degree $m$ and an inner
polynomial of degree $d$ has the same $O(n^2)$ dense bound: its successive
multiplications have costs $O(jd^2)$, and
$\sum_{j\leq m}jd^2=O(m^2d^2)$.

There are at most $\tau(n)-2$ proper degree divisors, where $\tau$ is
the divisor-counting function. Testing every one, or stopping at a first
success, therefore costs at most $O(n^2\tau(n))$ field operations.
Early rejection often makes individual unsuccessful trials cheaper.

\subsection{Cost of one complete recursive decomposition}
At a split of degree $N=ab$ with $a,b\geq2$,
\[
 a^2+b^2\leq\frac{N^2}{2}.
\]
For example, dividing by $a^2b^2$ reduces this to
$1/a^2+1/b^2\leq1/2$. Thus the sum of squared node degrees at successive
levels of the recursion tree decreases geometrically. The sum over all
nodes is at most $2n^2$. Every node degree divides $n$, so its divisor
count is at most $\tau(n)$.

It follows that the entire one-chain decomposition still has the field
operation bound
\begin{equation}\label{eq:complexity}
 O(n^2\tau(n)),
\end{equation}
including the final dense recomposition. A trial needs $O(n)$ coefficient
slots apart from the representation sizes of coefficients; retaining all
reports or memoized enumeration states needs more storage.

This counts coefficient-field operations, not expression parsing or
integer-divisor generation. It is an elementary bound for this
implementation, not a claim of optimal polynomial-decomposition complexity. The classical literature
contains more sophisticated algorithms \cite{kozen-landau}. The supplied
code favors transparent exact operations over asymptotic optimization.

\subsection{Bit cost and coefficient-field growth}
A field-operation count is not a bit-complexity bound. Let
$D=[K:\Q]$. Exact arithmetic in $K$ depends on $D$, rational coefficient
heights, the chosen representation, and the sizes of intermediate
numerators and denominators. Independent radicals can generate a field
whose degree grows rapidly with the number of generators. Canonical
\wt{RootReduce} conversions can dominate the polynomial-vector work.

The implementation is dense. A sparse expression such as $x^{1000000}$
is expanded into a coefficient list of length $1000001$; no sparse
complexity claim is made. The enumeration cost additionally depends on
the number of chains and the memoized states explored. A finite output
cap is not an asymptotic bound on that search.

\section{Worked examples}
\label{sec:examples}
The expressions in this section are exact mathematical results, checked
by the independent reference implementation. Wolfram snippets specify
usage and expected values, not a transcript of an executed native kernel.

\subsection{The radical-coefficient example from the question}
Put $a=\sqrt2$ and
\begin{align*}
 p(x)={}&3+3a+(14+4a)x+(12+26a)x^2\\
       &+(56+8a)x^3+(8+48a)x^4+48x^5+16a x^6.
\end{align*}
The only nontrivial right degrees to test are two and three.
For $d=2$, $m=3$ and
\[
 c_1=\frac{48}{16\sqrt2}=\frac{3}{\sqrt2},
 \qquad b_1=\frac{c_1}{3}=\frac1{\sqrt2}.
\]
Thus the unique candidate is
\[
 h(x)=x^2+\frac{x}{\sqrt2}.
\]
Repeated division gives constant digits, and the outer polynomial is
\begin{equation}\label{eq:originalF}
 F(y)=3+3\sqrt2+(8+14\sqrt2)y
       +(8+24\sqrt2)y^2+16\sqrt2\,y^3.
\end{equation}
Both component degrees are prime, so the pair is complete.

The original normalization in the question uses $g=\sqrt2\,h=x+\sqrt2 x^2$.
Its outer component is $F(y/\sqrt2)$, namely
\[
 3+3\sqrt2+(14+4\sqrt2)y+(4+12\sqrt2)y^2+8y^3.
\]
The package's different printed form is therefore the same decomposition
with a canonical inner normalization.

For $d=3$, $m=2$, the candidate is
\[
 h_3=x^3+\frac{3\sqrt2}{4}x^2+
          \left(\frac{15}{16}+\frac{\sqrt2}{8}\right)x.
\]
Already the first remainder of $p$ modulo $h_3$ is
\begin{equation}\label{eq:obstructionoriginal}
 \left(\frac32+\frac{51\sqrt2}{16}\right)x^2+
 \left(\frac{51}{16}+\frac{3\sqrt2}{4}\right)x+3+3\sqrt2.
\end{equation}
It is nonconstant, so degree three is conclusively excluded. There is
exactly one normalized complete chain for this example.

\begin{lstlisting}
Clear[x];
p = 3 + 3 Sqrt[2] + (14 + 4 Sqrt[2]) x +
    (12 + 26 Sqrt[2]) x^2 + (56 + 8 Sqrt[2]) x^3 +
    (8 + 48 Sqrt[2]) x^4 + 48 x^5 + 16 Sqrt[2] x^6;
chain = AlgebraicDecompose[p, x];
AlgebraicVerifyDecomposition[p, chain, x,
  "RequireComplete" -> True, "RequireNormalized" -> True]
(* Expected: True *)
AlgebraicDecomposeAtDegree[p, x, 3]
(* Expected: Missing["NotDecomposableAtDegree", 3] *)
\end{lstlisting}

\subsection{The nested degree-24 example}\label{sec:nested}
Set $p_{24}(x)=p(x^4-x+1)$. A normalized complete chain is
\begin{align*}
 A(y)={}&141+105\sqrt2+(168+142\sqrt2)y\\
       &+(56+72\sqrt2)y^2+16\sqrt2\,y^3,\\
 B(y)={}&y^2+\left(2+\frac1{\sqrt2}\right)y,\\
 C(x)={}&x^4-x.
\end{align*}
Thus $p_{24}=A\circ B\circ C$ and the degree sequence is $(3,2,4)$.
One can derive the coefficients without expanding a degree-24 polynomial:
write $x^4-x+1=C+1$, expand $h(C+1)$ as a quadratic in $C$, move its
constant into the cubic outer component, and use~\eqref{eq:originalF}.

The quartic $C$ is indecomposable. Its only possible nontrivial inner
degree is two; the top coefficients give candidate $x^2$, and its first
remainder is $-x$. Both ascending and descending trial orders must still
split any decomposable factor produced earlier. Processing both recursive
children makes this independent of the initially selected degree.

\begin{lstlisting}
p24 = Expand[p /. x -> x^4 - x + 1];
c24 = AlgebraicDecompose[p24, x];
Exponent[#, x] & /@ c24
(* Expected: {3, 2, 4} *)
AlgebraicVerifyDecomposition[p24, c24, x,
  "RequireComplete" -> True, "RequireNormalized" -> True]
(* Expected: True *)
\end{lstlisting}

\subsection{A quintic \texttt{Root} coefficient}
Let $\alpha$ be the first Wolfram polynomial root of $z^5-z-1$ and put
\[
 h=x^2+\alpha x,\qquad f=y^3+\alpha y+1,
 \qquad p=f(h).
\]
The exact normalized complete chain is $(f,h)$. The algorithm uses
$\alpha^5=\alpha+1$ in exact coefficient arithmetic and does not need a
radical expression for $\alpha$.
\begin{lstlisting}
alpha = Root[#^5 - # - 1 &, 1];
pRoot = Expand[(x^2 + alpha x)^3 +
  alpha (x^2 + alpha x) + 1];
AlgebraicDecompose[pRoot, x]
(* Expected exact values:
   {x^3 + alpha x + 1, x^2 + alpha x} *)
\end{lstlisting}
The reference tests also use randomized decompositions in this quintic
field. The prepared native suite includes a selected complex conjugate
as well as the first root, so that value selection is not replaced by a
minimal-polynomial-only comparison.

\subsection{Mixed radicals and complex algebraic values}
In $K=\Q(\sqrt2,\sqrt3,i)$, take
\[
 h=x^2+(\sqrt2+i)x,\qquad f=y^3+\sqrt3 y+1.
\]
The field has degree eight over $\Q$; the reference tests represent it
as an exact algebraic field and recover these prescribed components from
their expanded composition. There is no special real-coefficient branch
in the decomposition procedure. The complex coefficient is simply an
exact element of $K$.

\subsection{An indecomposable polynomial of composite degree}
For $p=x^6+x$, the degree-two candidate is $x^2$, and the degree-three
candidate is $x^3$. Each first remainder is $x$. Therefore neither
candidate succeeds, and the complete chain is the singleton $\{x^6+x\}$.
This example tests the difference between a prime-degree shortcut and an
actual negative decision at composite degree.

A report contains the two obstructions:
\begin{lstlisting}
report = AlgebraicDecompositionReport[x^6 + x, x];
report["Status"]
(* Expected: "Indecomposable" *)
report["DecisionComplete"]
(* Expected: True *)
\end{lstlisting}

\subsection{Two complete decompositions of a Chebyshev polynomial}
For
\[
 T_6(x)=32x^6-48x^4+18x^2-1,
\]
the two normalized complete chains are
\begin{align*}
 &(32y^3-48y^2+18y-1)\circ x^2,\\
 &(32y^2-1)\circ\left(x^3-\frac34x\right).
\end{align*}
Their ordered degrees are $(3,2)$ and $(2,3)$. Direct substitution
verifies both identities. Since two and three are the only proper degree
divisors, fixed-degree uniqueness also proves that this list is complete.

\begin{lstlisting}
all = AlgebraicDecomposeAll[ChebyshevT[6, x], x];
all["ReturnedCount"]
(* Expected: 2 *)
all["EnumerationComplete"]
(* Expected: True *)
AlgebraicDecomposeAll[x^12, x, "MaxDecompositions" -> 2]
(* Three chains exist; returns two with EnumerationComplete -> False. *)
AlgebraicDecomposeAll[x^12, x, "MaxDecompositions" -> 3]
(* Returns all three with EnumerationComplete -> True. *)
\end{lstlisting}

\section{Audit of the derivative-factor approach in the question}
\label{sec:audit}
The chain-rule method is mathematically legitimate: if $p=f\circ g$,
then $g'$ divides $p'$. In characteristic zero, integrating a nonzero
scalar multiple of $g'$ recovers $g$ up to precisely the affine freedom
that normalization removes. An exhaustive search through divisors,
followed by an exact composition test, can therefore work. The difficulty
is not the existence of an algorithm; it is the avoidable combinatorial
search and some correctness details in the supplied implementation.

\subsection{A single distinct derivative factor is not a rejection test}
The code returns immediately when the nonconstant part of \wt{FactorList}
has fewer than two entries. Factor-list entries count \emph{distinct}
factors, not their multiplicities. For example,
\[
 p=x^4,
 \qquad p'=4x^3,
 \qquad p=x^2\circ x^2.
\]
There is only one distinct nonconstant derivative factor, but a valid
nontrivial decomposition exists. The early return misses it. Multiplicity
must be considered if this method is retained.

\subsection{Subset enumeration duplicates divisors}
Suppose the factorization is $p'=c\prod_i q_i^{e_i}$. The multiset of
individual factor copies has $E=\sum_i e_i$ members and $2^E$ subsets.
Many subsets give the same divisor when multiplicities occur. The number
of distinct monic divisors is only $\prod_i(e_i+1)$.
Exponent-vector enumeration avoids those duplicates, but the resulting
search can still be very large. In contrast, the factorization-free
method has only one candidate for each proper degree divisor of $n$.

The derivative degree filters should also use the exact constraints
\[
 d=\deg g=1+\deg g',\quad d\mid n,\quad 2\leq d\leq n/2.
\]
In particular, the upper bound for $\deg g'$ is $n/2-1$, not $n/2$.
Constructing an outer polynomial before checking $d\mid n$ leads to
needless or malformed candidate systems based on the noninteger quotient
$n/d$.

\subsection{The subtle issue of recursing only on the outer factor}
It would be too strong to say that outer-only recursion is always wrong
for the exact exhaustive subset order in the supplied code. Suppose the
selected $g$ were itself decomposable, $g=A\circ B$. Then
$g'=(A'\circ B)B'$, and $B'$ is a proper divisor of $g'$ with fewer
irreducible factor copies. Its integrated candidate also gives a right
component of $p$. If every subset is tested in increasing cardinality
with a complete exact composition test, the candidate from $B'$ would
have succeeded earlier. Under those assumptions, the first successful
$g$ is indecomposable.

That is an \emph{order-dependent proof obligation}, not an automatic
property of every derivative search. A refactoring to degree heuristics,
a different divisor generator, or an early search limit can invalidate
it. The new package recursively decomposes both components explicitly,
so it requires no such implicit guarantee. Ascending degree order gives
an additional simple justification, but descending order remains correct.

\subsection{The power-merging rewrite changes composition incorrectly}
The supplied rewrite replaces adjacent $x^n$ and $x^m$ by $x^{n+m}$.
For composition, however,
\[
 x^n\circ x^m=x^{nm},\qquad\text{not }x^{n+m}.
\]
Thus whenever the rule applies with $nm\ne n+m$, it changes the
represented polynomial. Replacing addition by multiplication restores
composition equality, but still merges factors and can destroy a
\emph{complete} decomposition. A complete-chain API should not merge
adjacent prime-degree powers at all.

\subsection{Exact algebraic identity checks should be explicit}
A generic symbolic zero heuristic is unnecessary for algebraic numbers.
Reduce each coefficient exactly and check equality to zero. The new
package also validates the coefficient domain before starting, so an
unsupported parameter, an inexact input, and a proved negative
decomposition result have different return values.

The derivative-factor strategy can be repaired, but no repair is needed
for the new core: it never factors $p'$ and never enumerates divisors of
that derivative. The only divisors it enumerates are integer divisors of
$\deg p$.

\addtocontents{toc}{\protect\newpage}
\section{Public Wolfram Language interface}
\label{sec:api}

\subsection{Loading and evaluation conventions}
The simplest installation is to load the package file directly:
\begin{lstlisting}
base = "C:/path/to/AlgebraicDecomposition";
Get[FileNameJoin[{base, "Kernel", "AlgebraicDecomposition.wl"}]];
Clear[x];
\end{lstlisting}
The package exports symbols in the context
\wt{AlgebraicDecomposition\textasciigrave}. It does not unprotect,
replace, or add definitions to the system \wt{Decompose}.
The supplied \wt{Kernel/init.m} and \wt{PacletInfo.wl} also describe
an optional package/paclet layout. No paclet installation is required
for direct \wt{Get}.

The intended platform is Wolfram Language 13.0 or later. That is a
compatibility target, not a measured minimum: no native kernel execution
was available for this release. Run the supplied MUnit suite in the target
kernel before relying on native compatibility.

\subsection{One chain and one requested degree}
\begin{lstlisting}
AlgebraicDecompose[p, x, "DegreeOrder" -> "Ascending"]
AlgebraicDecompose[p, x, "DegreeOrder" -> "Descending"]
AlgebraicDecomposeAtDegree[p, x, d]
\end{lstlisting}
The first two calls select the order of tested right degrees and return
one complete normalized chain. Different orders may return different
valid chains; they need not do so for every polynomial.

A successful fixed-degree call returns $\{f,h\}$. When $d$ is a valid
proper divisor but its unique candidate fails, the result is
\begin{lstlisting}
Missing["NotDecomposableAtDegree", d]
\end{lstlisting}
An invalid requested degree, such as $d=1$, $d=n$, or $d\nmid n$,
returns a \wt{Failure}, not \wt{Missing}. This distinction separates an
invalid question from a mathematically negative answer to a valid one.

\subsection{All right components and all complete chains}
\begin{lstlisting}
AlgebraicRightComponents[p, x]
AlgebraicDecomposeAll[p, x, "MaxDecompositions" -> Infinity]
\end{lstlisting}
The right-component call returns all nontrivial normalized pairs
$\{f,h\}$, in increasing order of $\deg h$. It returns an empty list
when there is no such pair. These are two-component decompositions, not
necessarily complete chains; either component may itself decompose.

The all-chain call returns an association:
\begin{lstlisting}
<|"Decompositions" -> {chain1, chain2, ...},
  "EnumerationComplete" -> True,
  "ReturnedCount" -> numberOfReturnedChains|>
\end{lstlisting}
The cap must be \wt{Infinity} or a positive integer. When truncation is
detected, \wt{EnumerationComplete} is \wt{False}. The reported count
is only the number of returned chains. For a constant, linear, or
indecomposable input, there is one terminal chain, the singleton
$\{p\}$, under the API convention.

\subsection{The decision report}
\begin{lstlisting}
AlgebraicDecompositionReport[p, x]
\end{lstlisting}
The top-level association has the keys \wt{Polynomial}, \wt{Degree},
\wt{Status}, \wt{DecisionComplete}, and \wt{Trials}. Status is one
of \wt{Constant}, \wt{Linear}, \wt{Decomposable}, or
\wt{Indecomposable}. The zero polynomial is reported with degree
\wt{-Infinity}. A normally returned report tests every proper divisor,
so \wt{DecisionComplete} is \wt{True}; an arithmetic failure produces
a \wt{Failure} instead of a misleading incomplete negative report.

Each trial contains the right degree, the outer degree, the normalized
candidate, and \wt{DecomposableAtDegree}. A success adds
\wt{OuterComponent}. A failure adds \wt{ObstructionDigit},
\wt{NonconstantRemainder}, and \wt{PreviousConstantDigits}. Digit
indices start at zero, as in~\eqref{eq:hadic}.

\subsection{Composition and verification}
\begin{lstlisting}
AlgebraicCompose[chain, x]
AlgebraicVerifyDecomposition[p, chain, x,
  "RequireComplete" -> False,
  "RequireNormalized" -> False]
\end{lstlisting}
Composition uses the outermost-first convention. In particular,
\wt{AlgebraicCompose[\{f,g,h\},x]} computes $f(g(h(x)))$.
The empty chain represents the identity $x$.

Verification always checks exact equality to $p$. With
\wt{RequireNormalized} set to \wt{True}, all components after the
first must be monic, zero at zero, and nonconstant. With
\wt{RequireComplete} set to \wt{True} for a nonlinear input, every
component must have degree at least two and have no nontrivial right
component. For a constant or linear input, completeness means a singleton
chain. Consequently the empty chain verifies the identity as a composition
but is not called a complete decomposition by this convention.

An equality or structural mismatch returns \wt{False}. An invalid
coefficient domain returns \wt{Failure}, not \wt{False}. Completeness
verification repeats exact decomposition tests on the components, so it
can cost substantially more than an identity-only check.

\subsection{Options and failures}
Options in this release use string names. Unknown options and values
outside the documented choices produce \wt{Failure} objects.
The principal tags are summarized below.
\begin{center}
\begin{tabular}{>{\raggedright\arraybackslash}p{0.33\linewidth}>{\raggedright\arraybackslash}p{0.57\linewidth}}
\toprule
\textbf{Tag} & \textbf{Meaning}\\
\midrule
\wt{InvalidArguments} & Wrong arity or argument kind.\\
\wt{InvalidVariable} & A numeric symbol was supplied as the variable.\\
\wt{NotPolynomial} & The evaluated expression is not polynomial in the variable.\\
\wt{InexactCoefficient} & An inexact real value remains in the coefficient list.\\
\wt{NonAlgebraicCoefficient} & Exact algebraic membership could not be certified.\\
\wt{InvalidRightDegree} & The requested degree is not a proper divisor greater than one.\\
\wt{UnknownOption}, \wt{InvalidOption} & Unsupported option name or value.\\
\wt{AlgebraicArithmetic} & Exact coefficient reduction did not complete normally.\\
\wt{InternalInvariant}, \wt{InternalVerification} & A required internal property or exact recomposition check failed.\\
\bottomrule
\end{tabular}
\end{center}
An unsupported computation is never intentionally converted to the
singleton $\{p\}$, because that would falsely suggest indecomposability.
The package has no probabilistic success mode or hidden numerical
fallback.

\section{Validation, reproducibility, and what remains untested}
\label{sec:validation}

\subsection{Executed exact reference checks}
The archive includes a separate Python implementation using exact SymPy
coefficient domains. The recorded run used Python 3.13.5 and SymPy 1.14.0
with deterministic seed 206618. It completed 574 assertions, all passing,
in approximately 5.654 seconds in the preparation environment. This is
an environment-specific check duration, not a Wolfram performance claim.
The individual results are in \wt{validation/reference\_results.json}.

The tests include the original radical example; the nested degree-24
example; monomial enumeration counts and cap boundaries; both Chebyshev
chains; composite-degree indecomposable examples; constants and linear
polynomials; 80 randomized constructed pairs over $\Q$; 45 over
$\Q(\sqrt2)$; a quintic number field and 12 additional randomized
pairs there; and the mixed complex degree-eight field.

Several checks deliberately use operations independent of the
decomposition code. Forty polynomial divisions are compared with
SymPy's division algorithm. Twenty candidate computations are compared
with a formal binomial-series expansion. Random constructed examples
supply known normalized right and outer factors through
\eqref{eq:normalization}, rather than asking another decomposition
routine to invent an expected answer. Exact recomposition is checked
throughout.

These checks support the mathematical recurrence, indexing, polynomial
arithmetic, recursive behavior, and expected results. They do not execute
Wolfram pattern matching, option dispatch, \wt{RootReduce}, or the native
input-domain checks.

\subsection{A useful warning from the external oracle}
The installed SymPy 1.14.0 decomposition routine was also tried as a
secondary comparator. It missed some known decompositions with explicit
field domains. A minimal observed example was
\begin{lstlisting}[language=Python]
import sympy as s
x = s.symbols('x')
p = (x**3 + x**2 + x)**2
s.Poly(p, x, domain=s.QQ).decompose()
# Observed in the installed 1.14.0: a singleton degree-six polynomial.
\end{lstlisting}
The input visibly supplies an outer square and a cubic inner factor.
Inspection of that installed implementation found a missing term in its
right-candidate recurrence, producing $x^3+x^2+3x/2$ instead of the
correct $x^3+x^2+x$ over $\Q$. Integer-domain quotient behavior could
accidentally hide this particular defect.

Eighteen of 24 selected constructed-composite comparisons disagreed in
degree multiset. These observations are retained in the JSON results,
not discarded. The archive's \wt{validation/ORACLE\_NOTE.md} gives the
minimal reproduction, the affected indexing, and the diagnosed update.
This is a statement about the installed version that was inspected, not
an assertion about all or current SymPy releases. No SymPy patch or
upstream submission is included.

Accordingly the external decomposition oracle is not used as the
completeness authority for the 574 passing checks. The known constructions,
formal-series checks, exact division checks, and the proofs above are
independent of that oracle's candidate routine. The observed discrepancy
also illustrates why recomposition alone cannot establish completeness:
returning the entire input as a singleton always recomposes correctly.

\subsection{Prepared native tests, not a native test run}
The package includes 107 MUnit \wt{VerificationTest} cases, using exact
comparisons. These cover input contracts, original and nested expected
chains, algebraic-number representations, selected root conjugates,
exact leading cancellation, principal radical semantics, rejection of
inexact and transcendental coefficients, unknown options, empty-chain
conventions, completeness verification, and enumeration limits. Forty
of the cases are frozen randomized rational or quadratic-field
compositions; the others include explicit analytic and negative fixtures.
Wolfram documents file-based execution through \wt{TestReport}
and \wt{VerificationTest} \cite{wl-testreport,wl-verificationtest}.

\begin{tcolorbox}
\textbf{Native execution was not available.} The Wolfram connector was
attempted for both context and evaluation and returned HTTP 404. No local
Wolfram kernel was available. Therefore the 107 MUnit tests are supplied
but unexecuted. They must not be added to the 574 executed reference
checks to produce a misleading total of native or executed tests.
\end{tcolorbox}

Six Wolfram source/test files passed a lexical check for balanced
brackets, nested comments, and strings. The same script checks the package
core for prohibited reliance on factoring, solving, built-in decomposition,
or numerical fallbacks. This is a lexical inspection, not a Wolfram parser
or a substitute for native execution.

\subsection{Reproducing the checks}
From the extracted archive root, run
\begin{lstlisting}[language={}]
python validation/run_reference_tests.py
python validation/wl_static_check.py
wolframscript -file Tests/RunTests.wls
\end{lstlisting}
The first two commands reproduce the executed checks. The third is the
native test step still to perform in an available kernel. It loads the
package, prints the kernel version, runs the test file, and exits nonzero
unless the report succeeds.

The runner prefers the modern \wt{ReportSucceeded} report property and
includes a fallback using older success/failure-count properties. That
compatibility code is itself unexecuted. Alternatively, after loading the
package in a fresh notebook kernel, run
\begin{lstlisting}
TestReport[FileNameJoin[
  {base, "Tests", "AlgebraicDecomposition.wlt"}]]
\end{lstlisting}
Use a fresh kernel to avoid definitions attached to symbols in the test
fixtures. The archive records native validation as pending rather than
claiming a passed platform matrix.

\section{Boundaries and further implementations}
\label{sec:extensions}

\subsection{Positive characteristic is a genuinely different problem}
The triangular uniqueness proof only needs the outer degree $m$ to be
invertible, but the convenient differential recurrence also divides by
$k$. In positive characteristic it cannot be transplanted blindly, even
in some cases where the triangular argument remains valid. Characteristic
zero avoids both issues.

When the outer degree is divisible by the characteristic, same-degree
uniqueness itself can fail. In $\mathbb F_4$, let $a$ run through the
three nonzero elements. Since $a^3=1$ and $2=0$,
\begin{align*}
 (x^2+a^2x)\circ(x^2+ax)
 &=(x^2+ax)^2+a^2(x^2+ax)\\
 &=x^4+a^2x^2+a^2x^2+a^3x\\
 &=x^4+x.
\end{align*}
The three different right components are all monic, zero-constant, and
of degree two. This directly disproves a characteristic-free version of
Theorem~\ref{thm:rigidity}. The supplied package intentionally has no
\wt{Modulus} option.

\subsection{Parameters and specialization}
The recurrence and monic division can be implemented over a rational
function field such as $\Q(u_1,\ldots,u_s)$. That would decide generic
decomposability in the parameters. Specialization is a separate issue:
a leading coefficient may vanish, an intermediate denominator may vanish,
or obstruction coefficients may become zero and create a new decomposition.

A parameter-aware extension should return conditions on parameter space,
not silently substitute symbolic equality heuristics for exact zero
testing. The closed conditions described in Section~\ref{sec:ritt}
give a starting point: clear valid denominators on each degree chart,
then analyze the simultaneous vanishing of all obstruction coefficients.
No such condition solver is included in this release.

\subsection{Faster arithmetic and candidate construction}
The most direct performance upgrade is an explicit common-number-field
backend, with coefficient arithmetic in a fixed quotient field. It can
reuse the minimal polynomial and rational basis representation instead
of asking \wt{RootReduce} to canonicalize unrelated intermediate syntax.
The decomposition core need not change if the backend preserves exact
zero and arithmetic contracts.

For large $d$, formal root extraction can use precision doubling. With
$B$ correct modulo $t^k$, a Newton correction for $B^m-C$ gives a root
modulo approximately $t^{2k}$:
\[
 B_{\rm new}=B-\frac{B^m-C}{mB^{m-1}}\pmod{t^{2k}}.
\]
The denominator is a unit of the formal series ring, because its constant
term is $m\ne0$. Fast multiplication, inversion, and truncated powers can
then replace the quadratic scalar recurrence. These are potential
optimizations, not implemented alternatives selected by a package option.

\subsection{Modular screening must respect specialization}
A future rational or number-field backend could reduce trial computations
modulo suitable primes or prime ideals to reject many degrees cheaply.
Such a rejection is valid only when reduction is defined for every
coefficient used and the necessary leading coefficients and integer
denominators remain nonzero. Bad reduction can destroy a genuine
characteristic-zero decomposition or change its degree, so a modular
failure without these conditions is not a certificate.

Conversely a successful modular decomposition is only a screening result:
an obstruction can vanish modulo one prime without vanishing in
characteristic zero. Lift and reconstruct candidates if useful, but
retain an exact characteristic-zero recomposition and decision step.
The delivered implementation takes this conservative exact route directly.

\subsection{Sparse inputs, rational functions, and approximate data}
Sparse polynomial data may merit sparse coefficient storage or direct
recognition of power and Chebyshev forms. The current dense vectors are
not designed to exploit a huge gap between degree and number of terms.

Rational functional decomposition requires handling numerator/denominator
cancellation and M\"obius rather than merely affine equivalence. The
polynomial top-coefficient proof and one-candidate formula are not an API
for that different problem. Similarly, a polynomial identity holding only
approximately requires a norm, perturbation model, and numerical
conditioning analysis. Exact algebraic input is a deliberate boundary
of this implementation, not a request to reinterpret decimals as exact
nearby numbers.

\section{Conclusions}
The obstruction to algebraic-coefficient decomposition is not a lack of
an algorithm, and it is not necessary to factor the derivative over an
ever larger extension. For each candidate right degree, normalization
and the leading coefficients determine one possible inner polynomial.
The formal root recurrence constructs it inside the original coefficient
field. Repeated monic division then provides either its unique outer
polynomial or a finite exact obstruction.

These local decisions are enough to obtain one complete chain, all
normalized right components, or every normalized complete chain. The
same framework handles radicals and selected algebraic \wt{Root}
values without branch guessing. Ritt theory explains the degree
invariants and the genuine nonuniqueness, while the implementation's
correctness follows from the explicit candidate and membership proofs.

The archive contains the complete Wolfram source, examples, a native
test suite ready to run, and an independently executed exact reference
implementation. The outstanding validation step is native Wolfram kernel
execution; it is recorded as such rather than obscured by the passing
reference checks.

\appendix
\section{Compact implementation map}
\label{app:map}
\begin{center}
\begin{tabular}{>{\raggedright\arraybackslash}p{0.36\linewidth}>{\raggedright\arraybackslash}p{0.54\linewidth}}
\toprule
\textbf{Internal routine} & \textbf{Responsibility}\\
\midrule
\wt{inputVector} & Validate evaluated input, canonicalize coefficients, remove leading zeros.\\
\wt{red}, \wt{normalVector} & Exact coefficient reduction and vector normalization.\\
\wt{vAdd}, \wt{vSubtract}, \wt{vMultiply} & Dense exact polynomial-vector arithmetic.\\
\wt{vCompose}, \wt{vComposeChain} & Horner composition and outermost-first chain evaluation.\\
\wt{vDivideMonic} & Explicit monic long division.\\
\wt{rightCandidate} & Recurrence~\eqref{eq:recurrence} and zero-constant normalization.\\
\wt{outerData} & $h$-adic membership test and first obstruction digit.\\
\wt{trialData} & Combine candidate, membership, and exact pair verification.\\
\wt{rightPairs} & Successful normalized pairs over all proper degree divisors.\\
\wt{completeVectorChain} & Recursive decomposition of both sides of a successful split.\\
\wt{verifyChain} & Exact recomposition identity check.\\
\bottomrule
\end{tabular}
\end{center}

The full source is \wt{Kernel/AlgebraicDecomposition.wl}. The two core
listings in Section~\ref{sec:implementation} are reproduced for explanation;
the package file is the authoritative implementation to load. No generated
binary, external service, or online connection is needed by the package
itself once a Wolfram kernel is available.

\section{A small manual correctness checklist}
For a new implementation of the same algorithm, the following invariants
are particularly effective at detecting subtle mistakes:
\begin{enumerate}
\item The candidate is monic and zero-constant, with degree exactly $d$.
Its top $d$ coefficients make $p-a_nh^m$ have degree at most $n-d$.
\item The recurrence has the coefficient $(m+1)j-mk$, not $mj-mk$
or another expression missing the $j$ contribution from $C'B$.
\item Monic division satisfies $p=qh+r$ exactly and $\deg r<d$.
The zero polynomial has a consistent representation.
\item Success means \emph{all} $h$-adic digits are constant, not merely
the first one. A first constant remainder does not prove composition.
\item Both sides of a successful split are processed when constructing
a complete chain, unless an explicit proof guarantees one side is
already indecomposable.
\item Enumerating all binary parenthesizations is not the same as
enumerating distinct chains. Use the rightmost indecomposable component
or an equally rigorous uniqueness strategy.
\item A cap reached exactly at the last answer is not truncation.
Look ahead for another distinct chain before setting an incomplete flag.
\item Exact recomposition establishes equality, but not indecomposability.
Negative degree tests or known prime degrees are additionally required.
\end{enumerate}

\section{Archive layout and rebuilding the article}
\begin{lstlisting}[language={}]
AlgebraicDecomposition/
  README.md
  LICENSE.txt
  PacletInfo.wl
  Kernel/
    AlgebraicDecomposition.wl
    init.m
  Examples/Examples.wl
  Tests/
    AlgebraicDecomposition.wlt
    RunTests.wls
  Reference/exact_decomposition.py
  article/
    article.tex
    article.pdf
    build.sh
  validation/
    STATUS.md
    ORACLE_NOTE.md
    run_reference_tests.py
    reference_results.json
    worked_examples.txt
    make_native_tests.py
    native_test_count.txt
    wl_static_check.py
    wl_static_results.json
\end{lstlisting}
The TeX source is self-contained and uses a traditional bibliography;
no external bibliography processor is required. With a suitable TeX Live
installation, build it from the \wt{article} directory using
\begin{lstlisting}[language={}]
latexmk -pdf -interaction=nonstopmode -halt-on-error article.tex
\end{lstlisting}
The preamble names the required font and layout packages. Their font
files are not distributed in the archive. Original code, tests, reference
implementation, and exposition are supplied under MIT-0; cited works
retain their own rights.

\begin{thebibliography}{99}
\addcontentsline{toc}{section}{References}

\bibitem{mse}
Vladimir Reshetnikov, \emph{Is it possible to make Decompose work with
coefficients containing radicals?} Mathematica Stack Exchange,
question 206618, with answer 206622 and discussion, 2019.
The question, answer code, and Daniel Lichtblau's historical comment were
also supplied in the task materials.
\url{https://mathematica.stackexchange.com/q/206618/7288};
\url{https://mathematica.stackexchange.com/a/206622/7288}.

\bibitem{kozen-landau}
Dexter Kozen and Susan Landau, \emph{Polynomial decomposition algorithms},
Journal of Symbolic Computation \textbf{7} (1989), 445--456.
Bibliographic data and author abstract available on Dexter Kozen's
publication page:
\url{https://www.cs.cornell.edu/~kozen/Papers/papers_by_year.htm}.
The present article derives its working recurrence and correctness
arguments explicitly rather than relying on unquoted implementation
details from this reference.

\bibitem{zieve-mueller}
Michael E. Zieve and Peter M\"uller, \emph{On Ritt's polynomial decomposition
theorems}, preprint dated 23 July 2009. In particular, Section 2,
Corollary 2.12, Theorem 2.17, and Section 5.2.
\url{https://websites.umich.edu/~zieve/papers/peter.pdf}.

\bibitem{wl-decompose}
Wolfram Research, \emph{Decompose}, Wolfram Language documentation.
\url{https://reference.wolfram.com/language/ref/Decompose.html}.

\bibitem{wl-rootreduce}
Wolfram Research, \emph{RootReduce}, Wolfram Language documentation.
\url{https://reference.wolfram.com/language/ref/RootReduce.html}.

\bibitem{wl-tonumberfield}
Wolfram Research, \emph{ToNumberField}, Wolfram Language documentation.
\url{https://reference.wolfram.com/language/ref/ToNumberField.html}.

\bibitem{wl-algebraics}
Wolfram Research, \emph{Algebraics}, Wolfram Language documentation.
\url{https://reference.wolfram.com/language/ref/Algebraics.html}.

\bibitem{wl-testreport}
Wolfram Research, \emph{TestReport} and \emph{TestReportObject},
Wolfram Language documentation.
\url{https://reference.wolfram.com/language/ref/TestReport.html};
\url{https://reference.wolfram.com/language/ref/TestReportObject.html}.

\bibitem{wl-verificationtest}
Wolfram Research, \emph{VerificationTest}, Wolfram Language documentation.
\url{https://reference.wolfram.com/language/ref/VerificationTest.html}.

\end{thebibliography}
\noindent\small Web references were consulted on 7 September 2026.
External publications and documentation are cited, not redistributed.
\end{document}
